\documentclass[11pt]{article}
\usepackage{fontspec}
\usepackage[french,english]{babel}
\usepackage{amsmath}
\usepackage{amssymb}
\usepackage{amsthm}
\usepackage{mathrsfs}
\usepackage{tikz-cd}
\usepackage{enumitem}
\usetikzlibrary{babel}

\newcommand{\printedpage}[1]{\marginpar{\footnotesize p.~#1}}
\newcommand{\msp}[1]{\marginpar{\footnotesize [#1]}}
\newcommand{\Mon}{\mathrm{Mon}}
\newcommand{\Ob}{\mathrm{Ob}}
\newcommand{\Fl}{\mathrm{Fl}}
\newcommand{\Ss}{\mathrm{Ss}}
\newcommand{\card}{\mathrm{card}}
\newcommand{\cercle}[1]{\tikz[baseline=(c.base)]\node[draw,circle,inner sep=1pt](c){#1};}

\begin{document}

\printedpage{49}
% head: A. Grothendieck : Dérivateurs, Ch. XV, Édition des Universités de Montpellier II et Paris VII
% foot: Version 5 novembre 2010

\textit{and}
\[
\text{\textit{there exists a small subset} } TC_0 \subset \Mon, \text{ \textit{such that} } \Mon \cap W = \widetilde{TC_0}\ . \tag*{(7.2.5)} % [sic] (7.2.5) -> numéro probablement autre (peut-être (7.2.3)), le même numéro (7.2.5) est repris plus bas pour une autre formule et précède (7.2.4)
\]

\textit{(More precisely, if} $\pi \geq \pi_0$ \textit{is a cardinal such that} $W$ \textit{is} $\pi$\textit{-accessible, we can take} $TC_0 = (\Mon \cap W)^{\pi}$ \textit{(or a} small \textit{subset of} $(\Mon \cap W)^{\pi}$ \textit{having the same isomorphism classes).) Thus} (th.~1 of §1) \textit{the couple} $(TC = \Mon \cap W,\ F = (TC)_*)$ \textit{is a closed Quillen couple, and there is even a} functorial \textit{factorization of every arrow} $f$ \textit{as}
\[
f \ = \ p\, i \ ; \qquad i \ \in \ TC \ , \qquad p \ \in \ F \ . \tag*{(7.2.4)}
\]
\textit{If we set}
\[
C \ = \ \Mon \ , \qquad TC \ = \ \Mon \cap W \ , \qquad F \ = \ (TC)_* \ , \qquad TF \ = \ C_* \ , \tag*{(7.2.5)}
\]
\textit{the triple} $(W, C, F)$ \textit{is a closed Quillen triple, if and only if we have the following additional conditions.}
\[
\left\{
\begin{array}{lll}
\text{(7.2.6)} & \cercle{g}\ \text{\textit{(on $\mathcal{M}$ only)}} & \parbox[t]{0.36\textwidth}{Mon \textit{stable under cobase change.}} \\[4ex]
\text{(7.2.7)} & \underbrace{TF}_{\overset{\text{def}}{=}\,(\Mon)_*} \ \subset \ W & \parbox[t]{0.36\textwidth}{i.e. \textit{if} d) \textit{every} $f : X \longrightarrow Y$ \textit{which makes} $X$ \textit{an injective object of} $\mathcal{M}/Y$\textit{, is in} $W$ \textit{(hence also in} $W_{\mathrm{univ}}$\textit{).}}
\end{array}
\right.
\]

\msp{page 90}

\textsc{Proof of} 7.1. We can give a somewhat more delicate statement, by supposing given from the outset, not a subset $M_0$ of $\Ob \mathcal{M}$, but a subset
\[
\Mon_0 \ \subset \ \Mon \ , \tag*{(7.1)}
\]
satisfying the following condition

\begin{quote}
For every monomorphism $i : A \longrightarrow B$ in $\mathcal{M}$, there exists $i' \in \Mon_0$, $i' : A' \longrightarrow B'$, and an arrow $i' \longrightarrow i$ in $\underline{\Fl}(\mathcal{M})$,
\[
\begin{tikzcd}[row sep=large, column sep=large]
A' \arrow[r, "{i'}"] \arrow[d] & B' \arrow[d] \\
A \arrow[r, "i"] & B \ ,
\end{tikzcd}
\tag*{(7.2)}
\]
such that the corresponding arrow
\[
C \ = \ A \amalg_{A'} B' \ \longrightarrow \ B
\]
is a monomorphism, and $B' \longrightarrow B$ does not factor through $A$.
\end{quote}

That is what was aimed at by condition (7.1.1) on $M_0$, which implies (7.2) on $\Mon_0$ defined in terms of $M_0$ as said in 1°.

\printedpage{50}
% head: A. Grothendieck : Dérivateurs, Ch. XV, Édition des Universités de Montpellier II et Paris VII
% foot: Version 5 novembre 2010

In what follows, we suppose given a subset $\Mon_0 \subset \Fl(\mathcal{M})$ satisfying (7.2). In practice, it will often be deduced from a subset $M_0$ of $\Ob \mathcal{M}$, generating for monomorphisms, as made explicit in 1°, but we do not need the fact that $i' \longrightarrow i$ is a monomorphism (\textit{i.e.} (7.1.1) is stronger than what we need). E.g. if $B' \in M_0$ and $B' \longrightarrow B$ is given, not factoring through $A$, we shall be able to take $A' = B' \times_B A$ and $i' : A' \longrightarrow B'$ (to ensure $i' \in \Mon_0$, we shall therefore have to suppose $M_0$ stable under passage to subobjects), and chances are

\msp{page 91}

that $C \longrightarrow B$ will then indeed be a monomorphism \ldots

1°) We must prove
\[
\Mon \ \subset \ \widetilde{\Mon_0}
\]
assuming (7.2) (and the stability of $\mathcal{M}$ under filtering $\varinjlim$). This will follow directly from the

\textbf{Lemma 7.3.} \textit{Let $\mathcal{M}$ be a category stable under filtering direct limits, and $C_0$, $C$ subsets of} $\Fl(\mathcal{M})$\textit{. We shall suppose}
\begin{enumerate}[label={}]
\item[\cercle{a}] \textit{$C$ made up of monomorphisms, and}
\item[\cercle{b}] \textit{that the following condition (generalizing} (7.2) \textit{above) is satisfied.}
\end{enumerate}

\begin{quote}
\textit{For every arrow $i : A \hookrightarrow B$ in $C$ such that $i$ is not an isomorphism, there exists $i' : A' \longrightarrow B'$ in $\widetilde{C_0}$, and an arrow $i' \longrightarrow i$ in} $\underline{\Fl}(\mathcal{M})$\textit{, such that in the diagram}
\[
\begin{tikzcd}[row sep=large, column sep=large]
A' \arrow[r, "{i'}"] \arrow[d, hook] & B' \arrow[d] \arrow[rd, dotted] & \\
A \arrow[r, "\alpha"] \arrow[rr, dotted, bend right, "i"'] & D \ = \ A \amalg_{A'} B' \arrow[r, "\beta"] & B \ ,
\end{tikzcd}
\tag*{(7.3.1)}
\]
\textit{the arrow $\beta$ is in $C$, and $B' \longrightarrow B$ does not factor through $A$ (or, what amounts to the same, the subobject of $B$ defined by $D$ and $\beta$ is not equal to the one defined by $A$ and $i$).}
\end{quote}

\textit{We suppose moreover}
\begin{enumerate}[label={}]
\item[\cercle{c}] \textit{$C$ stable under filtering} $\varinjlim$\textit{, and}
\item[\cercle{d}] \textit{that for every $B \in \Ob \mathcal{M}$, the set of subobjects of $B$ is small.} (A most innocuous condition! This is a condition \textit{on} $\mathcal{M}$, not involving $C$).
\end{enumerate}
\textit{Then we have}
\[
C \ \subset \ \widetilde{C_0} \ , \qquad \text{\textit{hence} } \widetilde{C} \ \subset \ \widetilde{C_0} \ , \tag*{(7.3.2)}
\]
\textit{hence, if $C_0 \subset C$, we conclude}

\msp{page 92}
\[
\widetilde{C} \ = \ \widetilde{C_0} \ . \tag*{(7.3.3)}
\]

\printedpage{51}
% head: A. Grothendieck : Dérivateurs, Ch. XV, Édition des Universités de Montpellier II et Paris VII
% foot: Version 5 novembre 2010

(This lemma immediately implies the first assertion of 7.1, 1°.) Let us prove this lemma, \textit{i.e.} that every $i \in C$ is in $\widetilde{C_0}$, by constructing by transfinite induction a strictly increasing sequence of subobjects of $B$,
\[
A_\alpha \ \overset{j_\alpha}{\hookrightarrow} \ B \ , \qquad \text{with the } j_\alpha \in C \ ,
\]
and with:

\cercle{1°} $A_0 = A$. We shall prove that $A \xrightarrow{\ i_\alpha \,=\, i_{\alpha,0}\ } A_\alpha$ in $\widetilde{C_0}$, and there will be a final ordinal $\alpha$ such that $A_\alpha = B$\,; we shall therefore conclude from this that $i = i_\alpha$ is in $\widetilde{C_0}$.

If $i$ is an isomorphism, we know that $i \in \widetilde{C_0}$, there is nothing to prove (the transfinite construction stops at stage 0). The two steps of the construction, after the initial step 1°, are:

\cercle{2°} Let $A_\alpha$ be constructed. If $A_\alpha = B$, the construction is stopped. If $A_\alpha \xrightarrow{j_\alpha} B$ is not an isomorphism, we apply (7.3.1) to $j_\alpha \in C$, to find a $D \overset{j}{\hookrightarrow} B$, thus $j \in C$ and $D \supsetneqq A_\alpha$. Moreover,
\[
i_{\alpha+1,\alpha} \ : \ A_\alpha \ \longrightarrow \ A_{\alpha+1} \ \overset{\text{def}}{=} \ D \qquad \text{is in } \widetilde{C_0} \ , \tag*{(7.3.4)}
\]
since obtained from $i' \in \widetilde{C_0}$ by cobase change.

\cercle{3°} If $\alpha$ is a limit ordinal, we set
\[
A_\alpha \ = \ \varinjlim_{\beta < \alpha} A_\beta \ , \tag*{(7.3.5)}
\]

\msp{page 93}

and the canonical arrow
\[
A_\alpha \ \xrightarrow{j_\alpha} \ B
\]
is in $C$, since the $j_\beta$ are and $C$ is stable under filtering direct limits.

We have thus completed the description of the transfinite systems. For every $\alpha$, we have
\[
\card \alpha \ \leq \ \card \Ss(B) \ ,
\]
since the map $\beta \longmapsto A_\beta$ from the set $I_\alpha = \{\beta \mid \beta < \alpha\}$ to $\Ss(B)$ is \textit{strictly} increasing. Thus there exists a cardinal $\alpha$ such that $\card \alpha \leq \card \Ss(B)$, such that $A_\alpha = B$. But then $i \in \widetilde{C_0}$, as a transfinite direct limit (or composite) of arrows belonging to $\widetilde{C_0}$. (The condition $\forall\ \beta$, $i_{\beta+1,\beta} \in \widetilde{C_0}$ suffices to imply that all the transition morphisms $i_{\beta',\beta}$ are in $\widetilde{C_0}$.)

Let us return to 7.1, 2°. It remains to prove that if Mon is stable under cobase change, then we have
\[
\Mon \ = \ \widetilde{\Mon_0} \quad .
\]

But since Mon is already stable under composition, and under arbitrary filtering direct limits (hence also under direct factors), it follows by corollary 2 of theorem 1, §1 (p. 11), that $\Mon = \widetilde{\Mon_0}$\ \ . The factorization statement is theorem 1 of §1.

\printedpage{52}
% head: A. Grothendieck : Dérivateurs, Ch. XV, Édition des Universités de Montpellier II et Paris VII
% foot: Version 5 novembre 2010

\msp{page 94}

Let us prove 7.1, 3°, \textit{i.e.} formula (7.1.6),
\[
W \cap \Mon \ \subset \ \widetilde{(W \cap \Mon)^{\pi}} \ .
\]

For this, we again apply lemma 7.3 to
\[
C_0 \ = \ (W \cap \Mon)^{\pi} \ , \qquad C \ = \ W \cap \Mon \ .
\]

The conditions (\textit{a}), (\textit{b}), (\textit{c}), (\textit{d}) of the lemma are satisfied. This is clear for all, except for (\textit{b}), which calls for a non-tautological verification. Let then
\[
f \ : \ A \hookrightarrow B \ , \qquad f \in \Mon \cap W,\ f \text{ not an isomorphism} \ .
\]

We regard $B$ as a direct limit\footnote[33]{\textit{Of a $\pi$-adapted direct system.}}
\[
B \ = \ \varinjlim_{I} B_i \ , \qquad I \ = \ \Ss_{\pi}(B) \ .
\]
For every $i$, we set
\[
A_i \ = \ A \cap B_i \ = \ A \times_B B_i \ ,
\]
and
\[
f_i \ : \ A_i \ \hookrightarrow \ B_i \ ,
\]
hence
\[
f \ = \ \varinjlim_{i \in I} f_i \ .
\]
We need to know that the direct system of arrows $f_i$ is $\pi$-adapted, which leads us to postulate the following additional exactness condition on $\mathcal{M}$:

\msp{page 95}
\[
e')\quad \text{the filtering } \varinjlim \text{ in } \mathcal{M} \text{ commute with base change}\ , \tag*{(7.4)}
\]

it even suffices to suppose it in the very special case we have considered here.

Since $W$ is $\pi$-accessible, and $(f_i)_{i \in I}$ is $\pi$-adapted, since $\varinjlim f_i \in W$, it follows that the set $I_W$ of $i \in I$ such that $f_i \in W$ is cofinal. Since $A \overset{f}{\hookrightarrow} B$ is not an isomorphism, there is therefore an $i \in I_W$ such that $B_i \hookrightarrow B$ does not factor through $A$. We then have the diagram
\[
\begin{tikzcd}[row sep=large, column sep=large]
A_i \arrow[r, hook, "f_i"] \arrow[d, hook] \arrow[rd, phantom, "\text{\footnotesize cocart.}"] & B_i \arrow[d, hook] \arrow[rd, dotted, hook, "\text{\footnotesize incl.}"] & \\
A \arrow[r, hook, "\alpha_i"] \arrow[rr, dotted, hook, bend right, "f"'] & C_i \arrow[r, "\beta_i"] & B \ ,
\end{tikzcd}
\]

\printedpage{53}
% head: A. Grothendieck : Dérivateurs, Ch. XV, Édition des Universités de Montpellier II et Paris VII
% foot: Version 5 novembre 2010
with $f_i \in (C \cap W)^{\pi} = C_0$, hence $\alpha_i \in \widetilde{C_0}$. We need to know that $\alpha_i \in W$, hence an additional condition on $W$:

\begin{quote}
If we have a cocartesian square
\[
\begin{tikzcd}[row sep=large, column sep=large]
A \arrow[r, hook, "i"] \arrow[d] & B \arrow[d] \\
A' \arrow[r, "{i'}"] & B' \ ,
\end{tikzcd}
\tag*{(7.5)}
\]
with $i \in \Mon \cap W$, then $i' \in W$\footnote[34]{\textbf{N.B.} This condition is satisfied.}.
\end{quote}

(I have strengthened (7.1.4) accordingly).

\msp{page 96}

Since by hypothesis $f = \beta_i \alpha_i$ is in $W$, and $\alpha_i \in W$, it follows by (7.1.4), (\textit{a}) that $\beta_i \in W$. We shall have to know moreover that $\beta_i$ is a monomorphism, hence $\beta_i \in \Mon \cap W = C$, which will finish establishing the missing condition (\textit{b}) of 7.3. This is the exactness condition (7.2.1), which I evidently have to add to the hypotheses in 7.1, 3°. Having done that, lemma 7.3 applies and establishes 7.1, 3°.

\textsc{Proof of} 7.2. Let us set
\[
C \ = \ \Mon \ , \qquad TC \ = \ C \cap W \ , \qquad TC_0 \ = \ TC^{\pi} \ ,
\]
hence by 7.1, 3°, we have
\[
\underbrace{TC_0}_{\substack{\text{essentially}\\ \text{small}}} \ \subset \ TC \ \subset \ \widetilde{TC_0} \ = \ \widetilde{TC} \ ,
\]
and by theorem 1, §1 (p. 10-12), cor. 2, to check that $TC = \widetilde{TC_0}$ (\textit{i.e.} $TC$ is Q-closed on the left), it amounts to the same to see that it is Q-stable on the left, \textit{i.e.} stable under composition, under cobase change, under strict ordinal $\varinjlim$, under direct factors. Now the last two are implied by the fact that $TC = C \cap W$ is stable under filtering $\varinjlim$, seeing that $W$ is, and $C$ also by condition \cercle{e} (p. 76) on $\mathcal{M}$. On the other hand, $C = \Mon$ and $W$ being stable under composition, likewise for $C \cap W$. So there remains cobase

\msp{page 97}

change, and this is ensured by 7.2.3, (\textit{c}) (which we already needed in 7.1, 3°, to establish $C \cap W \subset \widetilde{TC_0}$). (\textbf{N.B.} the hypotheses of 7.2 are practically those of 7.1, 3°, except that we suppose that $W$ is a localizer (understood: mildly saturated), to be precise we have used, up to now, as an additional hypothesis, only the stability of $W$ under composition.)

It remains to establish that $(W, C = \Mon, F = (TC)_*)$ is a closed Quillen triple, if and only if the conditions (7.2.6) and (7.2.7), \textit{i.e.}
\[
\left\{
\begin{aligned}
& C = \Mon \ \ \text{stable under cobase change,} \\
& TF = C_* = (\Mon)_* \subset W,
\end{aligned}
\right.
\]

\end{document}
